\begin{description}
\item[(a)] From the formula $Q = \lambda S \Delta T t / d$, calculate the power
  $P$ of the heat flowing through a unit of the surface of ice:
  \[
    P = Q/t = \lambda S \Delta T / d.
  \]
  \hfill(1 point)

  Approximating the ice crust of Europa as a `wall', i.e.\ that the upper and
  lower surfaces are of equal area $S$, we obtain the power per unit surface
  area:
  \[
    P/S = \lambda \Delta T / d.
  \]
  \hfill(1 point)

  Substituting the data we obtain
  \[
    P/S = 86.5 \times 10^{-3}\ \mathrm{W\,m^{-2}}
        \approx 87\ \mathrm{mW\,m^{-2}},
  \]
  \hfill(1 point)

  similar to the value given for the Earth.

  The total power emitted inside Europa is therefore equal to
  \[
    4\pi R_{\mathrm{Eu}}^2 \, (P/S) = 2.65 \times 10^{12}\ \mathrm{W}.
  \]
  \hfill(2 points)

\item[(b)] The total power emitted inside the Earth is equal to
  \[
    4\pi R_{\oplus}^2 \times 70 \times 10^{-3}\ \mathrm{W\,m^{-2}}
      = 36 \times 10^{12}\ \mathrm{W},
  \]
  \hfill(1 point)

  which is about $13.5\times$ larger than the value for Europa. However, the
  mass ratio is 124. For the heat in Europa's interior to have a purely
  radioactive origin, the matter on Europa would have to contain an order of
  magnitude more radioactive elements per unit mass, which excludes this
  explanation.
  \hfill(3 points)

  Thus the answer must be tidal forces.
  \hfill(1 point)
\end{description}
